% Generated draft: requires mathematical and source review.
\documentclass[11pt]{article}
\usepackage[margin=25mm]{geometry}
\usepackage{fontspec}
\setmainfont{Latin Modern Roman}
\usepackage{amsmath,amssymb,mathtools,unicode-math}
\setmathfont{Latin Modern Math}
\usepackage{enumitem,needspace,xurl,hyperref}
\setlength{\parindent}{0pt}
\setlength{\parskip}{5pt}
\setlength{\emergencystretch}{3em}
\newcommand{\N}{\mathbb{N}}
\newcommand{\Z}{\mathbb{Z}}
\newcommand{\Q}{\mathbb{Q}}
\newcommand{\R}{\mathbb{R}}
\newcommand{\C}{\mathbb{C}}
\newcommand{\F}{\mathbb{F}}
\newcommand{\A}{\mathbb{A}}
\newcommand{\PP}{\mathbb{P}}
\newcommand{\E}{\mathbb{E}}
\newcommand{\eps}{\varepsilon}
\newcommand{\dd}{\,\mathrm{d}}
\begin{document}
\section*{UnsolvedMath ID: 30003008}
\textbf{Source problem code:} catalogue entry
\subsection{Definitions and mathematical statement}
The supplied record concerns a geometric mean \(G\) of convex bodies and operations \(M_{\mathrm{ag}}\) and \(M_{\mathrm{gh}}\), described respectively as arithmetic--geometric and geometric--harmonic iterations. It does not supply definitions of these operations, their domains, their iteration rules, or their convergence and limit conventions. Consequently, the following records the two intended questions schematically; a complete mathematical formulation requires the missing specifications.

Let \(\mathcal D\) denote the intended, unspecified class of convex bodies. The ambient space and dimension, admissible degeneracies, and any symmetry or origin conditions are not specified. The phrase ``natural class'' is also left undefined.

\begin{enumerate}
\item Does there exist a logarithmic map \(L\) on a suitable natural domain \(\mathcal D\) such that
\[
L\bigl(G(A,B)\bigr)=\frac{L(A)+L(B)}{2}
\qquad\text{for every }A,B\in\mathcal D?
\]
For this expression to be meaningful, \(G(A,B)\) must belong to the domain of \(L\), and the codomain of \(L\) must support the displayed addition and division by \(2\). Neither the codomain nor the defining requirements for a ``logarithmic map'' are supplied. In particular, the source does not state injectivity, normalization, regularity, or any other condition that would specify the intended meaning beyond the displayed equation.
\item Is the identity
\[
G\bigl(M_{\mathrm{ag}}(A,B),M_{\mathrm{gh}}(A,B)\bigr)=G(A,B)
\]
valid for every admissible pair \(A,B\) in the intended domain? Here the notation \(M_{\mathrm{ag}}(A,B)\) and \(M_{\mathrm{gh}}(A,B)\) denotes the outputs intended by the source's arithmetic--geometric and geometric--harmonic iterations. The record does not specify whether these outputs are limits, how the iterations are initialized or updated, or under which hypotheses the outputs exist. The outer application of \(G\) also requires these outputs to lie in its domain.
\end{enumerate}

These unresolved qualifications prevent either question from being a fully specified mathematical statement on the supplied evidence alone.
\subsection{Short English statement}
The source asks whether a logarithmic transformation converts a convex-body geometric mean into an arithmetic average, and whether taking the geometric mean of the arithmetic--geometric and geometric--harmonic iteration outputs recovers the original geometric mean. The domains, operations, iteration conventions, and meaning of ``logarithmic'' require clarification.
\subsection{Sources}
\begin{itemize}
\item Local UnsolvedMath registry: supplied record with permanent ID \(30003008\), source problem code \texttt{catalogue entry}, in the Oberwolfach Reports Open Problems collection. Project registry paths are \texttt{data/unsolvedmath/problems.json} and \texttt{catalog/problems.json}; the supplied material does not identify which file contains this entry.
\item Supplied report citation: \emph{Convex Geometry and its Applications} (2016), Oberwolfach Reports, pp.~3179--3269, DOI \texttt{10.4171/owr/2015/56}. Supplied source URL and literature-source link, labeled ``Convex Geometry and its Applications report'': \url{https://doi.org/10.4171/owr/2015/56}. No individual proposer is identified; the proposed year is recorded as 2015.
\item Supplied literature assessment, checked on \textbf{2026-08-21}: the two questions are described as open, and the curation reports that no general proof or counterexample was located. This is supplied dated curation derived from the Oberwolfach material, not an independently verified literature assessment. The supplied background requests independent verification of time-sensitive status and cited claims.
\item The registry labels the statement \texttt{corrected\_verified} and reports that the official report passage was checked directly and its question, definitions, and notation restored without strengthening the claim. This is a supplied verification claim; the present draft has not independently checked that passage, and the supplied record does not reproduce the needed definitions.
\end{itemize}
\end{document}
